\documentclass[11pt]{article}

% ---- Packages ----
\usepackage[utf8]{inputenc}
\usepackage[T1]{fontenc}
\usepackage{amsmath,amssymb,amsthm}
\usepackage{booktabs}
\usepackage{algorithm}
\usepackage{algorithmic}
\usepackage{hyperref}
\usepackage[capitalise]{cleveref}
\usepackage{microtype}
\usepackage{geometry}
\geometry{margin=1in}

% ---- Theorem environments ----
\newtheorem{theorem}{Theorem}[section]
\newtheorem{lemma}[theorem]{Lemma}
\newtheorem{proposition}[theorem]{Proposition}
\newtheorem{conjecture}[theorem]{Conjecture}
\theoremstyle{definition}
\newtheorem{definition}[theorem]{Definition}
\theoremstyle{remark}
\newtheorem{remark}[theorem]{Remark}

% ---- Custom commands ----
\newcommand{\LIS}{\operatorname{LIS}}
\newcommand{\LDS}{\operatorname{LDS}}
\newcommand{\Coll}{\operatorname{C}}

% ======================================================================
\begin{document}

\title{Longest Increasing Subsequences of Collatz Trajectories:\\
  A Null Model Decomposition}

\author{Mark T.\ Short\\
  \texttt{SevenevesAI}\\
  \texttt{mshort3@gmail.com}}

\date{\today}

\maketitle

% ======================================================================
\begin{abstract}
We study the longest increasing subsequence (LIS) of Collatz trajectories
as a quantitative probe of pseudorandomness. For starting values
$n = 1, \ldots, N$, the mean ratio of LIS length to $2\sqrt{k}$, where
$k$ is the trajectory length, is $0.567$ at $N = 10^6$ ---
significantly below the $0.812$ observed for random sequences of
matched lengths. Shuffling the trajectory values restores the ratio to
$0.813$, confirming that the deviation reflects temporal ordering, not the
value distribution. We decompose the deviation through a hierarchy of null
models: the $31.6\%$ up-step fraction accounts for most of the gap, while
the long descending suffix to~$1$ further depresses the ratio. The
residual beyond matched stochastic models is approximately $0.05$, small
but reproducible. The ratio varies across the generalized $3n+b$ family
($0.49$ to $0.61$), suggesting it encodes dynamical information about the
map. The asymptotic behavior of the ratio as $N \to \infty$ remains open.
Python code reproducing all experiments is available at
\url{https://github.com/sevenevesai/CollatzLIS}.
\end{abstract}

\noindent\textbf{Keywords:} Collatz conjecture, longest increasing
subsequence, pseudorandomness, trajectory geometry, null model hierarchy

\noindent\textbf{MSC 2020:} 11B83, 05A16, 37A45, 60C05

% ======================================================================
\section{Introduction}\label{sec:intro}

% --- 1.1 Context ---
\subsection{The Collatz map and trajectory statistics}

The Collatz map $T \colon \mathbb{Z}^+ \to \mathbb{Z}^+$ is defined by
\begin{equation}\label{eq:collatz}
  T(n) = \begin{cases} 3n+1 & \text{if } n \text{ is odd},\\
                        n/2  & \text{if } n \text{ is even}. \end{cases}
\end{equation}
The Collatz conjecture asserts that for every positive integer~$n$, repeated
application of~$T$ eventually reaches~$1$. Despite substantial progress ---
notably Tao's result~\cite{Tao2019} that almost all orbits attain almost
bounded values --- the conjecture remains open.

A complementary line of inquiry asks: \emph{to what extent do Collatz
trajectories resemble random sequences?} Statistical properties of
trajectories have been studied extensively, including stopping
times~\cite{Terras1976}, digit distributions~\cite{KontorovichMiller2005},
and density of iterates. The present paper contributes to this line by
studying the \emph{subsequence structure} of Collatz trajectories through
the lens of the longest increasing subsequence.

% --- 1.2 LIS as a probe ---
\subsection{LIS and the BDJ normalization}

Given a sequence $x_1, x_2, \ldots, x_k$ of real numbers, the
\emph{longest increasing subsequence} $\LIS(x)$ is the maximum length of
a strictly increasing subsequence. The celebrated theorem of Baik, Deift,
and Johansson~\cite{BDJ1999} establishes that for a uniformly random
permutation of $\{1, \ldots, k\}$,
\begin{equation}\label{eq:bdj}
  \frac{\LIS(\sigma) - 2\sqrt{k}}{k^{1/6}} \xrightarrow{d}
    F_2 \quad \text{as } k \to \infty,
\end{equation}
where $F_2$ is the Tracy--Widom distribution~\cite{TracyWidom1994}.

We emphasize that Collatz trajectories are \emph{not} random permutations:
they are deterministic sequences with values in $\mathbb{Z}^+$, exhibiting
strong directional bias and temporal correlation. The
quantity~$2\sqrt{k}$ serves here as a convenient normalization, not as a
theoretical prediction. Our comparisons are against empirical baselines
(random sequences, shuffled trajectories, matched stochastic models), not
against the BDJ asymptotics directly.

For a starting value $n \in \{1, \ldots, N\}$, let
$\Coll(n) = (n, T(n), T^2(n), \ldots, 1)$ denote the Collatz trajectory,
$k(n)$ its length, and $L(n) = \LIS(\Coll(n))$ the LIS of the trajectory.
The \emph{LIS ratio} is
\begin{equation}\label{eq:ratio}
  r(n) = \frac{L(n)}{2\sqrt{k(n)}},
\end{equation}
and the mean ratio over $\{1, \ldots, N\}$ is $\bar{r}(N) = N^{-1}
\sum_{n=1}^{N} r(n)$.

% --- 1.3 Main observation ---
\subsection{Main observation}

The central empirical finding is:

\begin{quote}
\emph{The mean LIS ratio $\bar{r}(N)$ for Collatz trajectories is
$0.567$ at $N = 10^6$, significantly below the random baseline of
$0.812$. The deviation is explained primarily by the $31.6\%$ up-step
fraction and the long descending suffix, with a residual of $\sim 0.05$
beyond matched stochastic models.}
\end{quote}

% --- 1.4 Contributions ---
\subsection{Contributions}

The contributions of this paper are:
\begin{enumerate}
  \item The empirical observation that the Collatz LIS ratio is
    significantly below random baselines, sustained across six orders of
    magnitude in~$N$ (\cref{sec:scaling}).
  \item A \emph{null model hierarchy} decomposing the deviation: random
    permutation $\to$ multiplicative random walk $\to$ Bernoulli process
    $\to$ Markov chain $\to$ Collatz, each adding one structural feature
    (\cref{sec:null}).
  \item A \emph{prefix--suffix decomposition} revealing that the low ratio
    is primarily a suffix effect: the prefix-only ratio is $0.85$, near the
    random baseline (\cref{sec:structure}).
  \item Evidence of \emph{parameter sensitivity}: the ratio varies across
    the $3n+b$ family ($0.49$ to $0.61$), serving as a dynamical
    fingerprint (\cref{sec:parameter}).
  \item An honest assessment of limitations: trajectory correlation,
    convergence uncertainty, and the absence of analytical explanation
    (\cref{sec:limitations}).
\end{enumerate}

\subsection{Outline}

The remainder of this paper is organized as follows.
\Cref{sec:setup} describes the experimental setup and LIS algorithm.
\Cref{sec:scaling} presents the scaling experiments.
\Cref{sec:null} develops the null model hierarchy.
\Cref{sec:structure} analyzes trajectory structure (directionality, peak
position, prefix--suffix decomposition).
\Cref{sec:parameter} reports the parameter sensitivity across the $3n+b$
family.
\Cref{sec:limitations} discusses trajectory correlation, convergence
uncertainty, and statistical caveats.
\Cref{sec:conclusion} summarizes and identifies open questions.

% ======================================================================
\section{Experimental Setup}\label{sec:setup}

% TODO: Python version, hardware, LIS algorithm, correctness verification,
%       reproducibility details

All experiments use Python~3.11. The LIS is computed via the patience
sorting algorithm in $O(k \log k)$ time. Correctness is verified by
exact match against brute-force enumeration on 1000 random sequences of
length~10 (0 mismatches out of 1000 trials; \cref{sec:correctness}).
Collatz trajectories are computed from starting values $n = 1, \ldots, N$
with no step limit; all trajectories in the standard $3n+1$ map reach~$1$.

Bootstrap confidence intervals use 1000 resamples. We note that trajectory
correlation (\cref{sec:correlation}) reduces the effective sample size to
approximately 500 at $N = 20{,}000$; reported CIs should be widened by a
factor of approximately~6.

Source code reproducing all experiments is available at
\url{https://github.com/sevenevesai/CollatzLIS}.

\subsection{LIS algorithm correctness}\label{sec:correctness}

The patience sorting implementation passes all 10 named edge cases
(empty sequence, single element, sorted, reverse-sorted, etc.) and
matches brute-force LIS on 1000 random sequences of length~10 with 0
mismatches. We consider the algorithm verified.

% ======================================================================
\section{Scaling Experiments}\label{sec:scaling}

\Cref{tab:scaling} reports the mean LIS ratio $\bar{r}(N)$ for
$N = 10^3$ to $N = 10^6$.

\begin{table}[ht]
\centering
\caption{Mean LIS ratio $\bar{r}(N)$ and inter-step change $\Delta$ as a
function of~$N$.}
\label{tab:scaling}
\begin{tabular}{@{}rrcc@{}}
\toprule
$N$ & $\bar{r}(N)$ & $\Delta$ & Random baseline \\
\midrule
1{,}000     & 0.632 & ---     & 0.81 \\
5{,}000     & 0.627 & $-0.005$  & 0.81 \\
10{,}000    & 0.621 & $-0.007$  & 0.81 \\
50{,}000    & 0.599 & $-0.022$  & 0.81 \\
100{,}000   & 0.590 & $-0.008$  & 0.81 \\
250{,}000   & 0.581 & $-0.010$  & 0.81 \\
500{,}000   & 0.574 & $-0.007$  & 0.81 \\
750{,}000   & 0.570 & $-0.004$  & 0.81 \\
1{,}000{,}000 & 0.567 & $-0.003$  & 0.81 \\
\bottomrule
\end{tabular}
\end{table}

The ratio is monotonically decreasing over the measured range. The
inter-step change $\Delta$ is diminishing but has not reached zero: at
$N = 10^6$ the ratio is still falling by $0.003$ per doubling step. The
ratio has \emph{not} converged at $N = 10^6$.

\paragraph{LIS/LDS asymmetry.}
The mean LIS is $13.0$ and the mean longest decreasing subsequence (LDS) is
$36.3$, giving a ratio $\LIS/\LDS = 0.36$. Collatz trajectories are almost
three times more ``decreasing'' than ``increasing,'' a direct consequence of
the $31.6\%$ up-step fraction.

% ======================================================================
\section{Null Model Hierarchy}\label{sec:null}

The deviation of the Collatz LIS ratio from the random baseline can be
decomposed by comparing against a hierarchy of stochastic models, each
capturing one additional structural property of Collatz trajectories.

\begin{table}[ht]
\centering
\caption{LIS ratio for Collatz and null models. Each model adds one
structural feature beyond the previous. The ``Gap'' column shows the
difference from the Collatz ratio.}
\label{tab:null}
\begin{tabular}{@{}llcc@{}}
\toprule
Model & What it captures & Ratio & Gap \\
\midrule
Random permutation       & Nothing (BDJ baseline)    & 0.812 & $+0.219$ \\
Shuffled Collatz values  & Value distribution only    & 0.813 & $+0.220$ \\
Multiplicative RW        & Step-size distribution     & 0.638 & $+0.045$ \\
Bernoulli ($p = 0.316$)  & Up-fraction only           & 0.571 & $-0.022$ \\
Markov (matched autocorr)& Up-fraction + alternation  & 0.478 & $-0.115$ \\
\midrule
\textbf{Collatz actual}  & \textbf{All structure}     & \textbf{0.593} & baseline \\
\bottomrule
\end{tabular}
\end{table}

\subsection{The shuffling test}\label{sec:shuffle}

The shuffled Collatz baseline ($0.813$) matches the random permutation
baseline ($0.812$) to within sampling noise. This confirms that the low
Collatz ratio reflects \emph{temporal ordering} rather than the marginal
value distribution. Destroying the temporal order restores the ratio to the
random baseline.

\subsection{The multiplicative random walk}\label{sec:multrw}

A multiplicative random walk with matched parameters --- $33\%$ up-steps of
magnitude $3\times$, $67\%$ down-steps of magnitude $0.5\times$ --- produces
ratio $0.638$, explaining most of the gap between Collatz ($0.593$) and
random ($0.812$). The remaining gap of $0.045$ is the residual attributable
to Collatz-specific structure beyond step-size statistics.

\subsection{The Bernoulli model}\label{sec:bernoulli}

An independent Bernoulli process with $p = 0.316$ (matching the Collatz
up-step fraction) produces ratio $0.571$, within $0.022$ of the Collatz
value. This is the key finding of the null model analysis: \emph{the
up-step fraction alone accounts for most of the LIS ratio deviation from
random}. No temporal correlation structure is needed to reproduce the
approximate ratio.

\subsection{The Markov model}\label{sec:markov}

The lag-1 autocorrelation of the Collatz up/down indicator sequence is
$-0.458$ (strongly alternating), reflecting the characteristic $3n+1 \to
/2 \to /2$ pattern. A Markov chain matching this autocorrelation gives
ratio $0.478$ --- \emph{lower} than the Collatz value, not higher. The
alternating structure does not explain the residual gap; in fact, it
overshoots. This suggests that the precise temporal arrangement of
increasing steps in Collatz trajectories is more favorable for LIS growth
than the strong alternation pattern alone would predict.

% ======================================================================
\section{Structural Analysis}\label{sec:structure}

\subsection{Directional bias}\label{sec:bias}

Only $31.6\%$ of transitions in Collatz trajectories are increasing. This
fraction is tightly concentrated across starting values, with the
distribution supported on $[0.2, 0.4]$ and no trajectory above $0.5$. The
asymmetry is a direct consequence of the map: a single $3n+1$ step (odd) is
typically followed by one or two $n/2$ steps (even), giving an up-step
fraction near $1/3$.

\subsection{Peak position}\label{sec:peak}

The average peak (maximum value) occurs at position $0.11$ --- just $11\%$
into the trajectory. A majority ($63\%$) of trajectories peak within the
first $10\%$ of their length. The longest contiguous increasing run occurs
at position $0.014$ on average, right at the start. Collatz trajectories
shoot up early, then cascade down to~$1$.

\subsection{Prefix--suffix decomposition}\label{sec:prefix}

This trajectory geometry has a direct consequence for LIS. We decompose
each trajectory into:
\begin{itemize}
  \item The \emph{prefix}: from $n$ to the peak value.
  \item The \emph{suffix}: from the peak value down to $1$.
\end{itemize}

The prefix-only LIS ratio is $0.853$, near the random baseline of $0.812$.
The full-trajectory ratio is $0.590$. The low ratio is driven primarily by
the long descending suffix: the suffix inflates the trajectory length $k$
(the denominator of $r$) while the LIS cannot grow during a monotone
descent (contributing at most $1$ to the numerator per descent segment).

This is a \emph{mechanical} consequence of trajectory shape, not a
number-theoretic signal. Any sequence that rises briefly, peaks early, and
then descends for most of its length will exhibit a depressed LIS ratio
relative to $2\sqrt{k}$.

% ======================================================================
\section{Parameter Sensitivity}\label{sec:parameter}

We extend the analysis to the generalized map $T_b(n) = 3n + b$ (odd) or
$n/2$ (even), for $b \in \{1, 3, 5, 7, 11\}$. Only $b = 1$ gives $100\%$
convergence to a cycle; other values have varying convergence rates under a
$10^5$-step guard.

\begin{table}[ht]
\centering
\caption{LIS ratio for the $3n+b$ family. Only converging trajectories are
included.}
\label{tab:parameter}
\begin{tabular}{@{}ccccc@{}}
\toprule
$b$ & Converge \% & Avg $k$ & Ratio & TW mean \\
\midrule
1   & 100.0 & 101.5 & 0.599 & $-3.59$ \\
3   & 0.03  &  90.7 & 0.612 & $-3.32$ \\
5   & 14.2  &  76.2 & 0.521 & $-3.88$ \\
7   & 69.2  &  80.9 & 0.489 & $-4.27$ \\
11  & 19.8  &  85.5 & 0.493 & $-4.30$ \\
\bottomrule
\end{tabular}
\end{table}

The ratio varies from $0.489$ ($b = 7$) to $0.612$ ($b = 3$), a spread of
$0.12$. This suggests that the LIS ratio encodes dynamical information
specific to each map variant, not merely a universal consequence of the
$3n+b$ structure. The standard $3n+1$ map falls in the middle of the range.

We note that the low convergence rate for $b \neq 1$ introduces selection
bias: only trajectories that reach a cycle are included, and these may not
be representative of typical orbits. The parameter sensitivity observation
should be treated as suggestive, not definitive.

% ======================================================================
\section{Limitations and Statistical Caveats}\label{sec:limitations}

\subsection{Trajectory correlation}\label{sec:correlation}

Collatz trajectories from nearby starting values share a large fraction of
their steps. On a sample of $20{,}000$ trajectories, $97.4\%$ of each
trajectory consists of a suffix shared with at least one previously
encountered trajectory, and $93\%$ of trajectories share more than $90\%$
of their length. The effective number of independent samples is
approximately $500$ out of $20{,}000$.

This means bootstrap confidence intervals computed naively (CI width
$0.003$ at $N = 10^6$) understate the true uncertainty by a factor of
roughly~$6$. The adjusted $95\%$ CI for $\bar{r}(10^6)$ is approximately
$[0.55, 0.59]$.

\subsection{Convergence uncertainty}\label{sec:convergence}

The ratio $\bar{r}(N)$ is still decreasing at $N = 10^6$. Three convergence
models give different asymptotic predictions:

\begin{table}[ht]
\centering
\caption{Convergence model fits and asymptotic predictions.}
\label{tab:convergence}
\begin{tabular}{@{}lccc@{}}
\toprule
Model & Asymptotic limit & RMSE & Predicted $\bar{r}(10^9)$ \\
\midrule
$r = a + b/\log N$  & 0.469 & 0.0023 & 0.536 \\
$r = a + b/\sqrt{N}$ & 0.572 & 0.0064 & 0.572 \\
$r = a + b/N^{0.25}$ & 0.549 & 0.0032 & $\approx 0.55$ \\
\bottomrule
\end{tabular}
\end{table}

The logarithmic model fits best (lowest RMSE) and predicts convergence to
$\sim 0.47$, but the power-law models predict limits near $0.55$--$0.57$.
We cannot determine the asymptotic limit with confidence from $N \leq 10^6$
data. A Rust or C implementation reaching $N = 10^8$ or beyond would
significantly reduce this uncertainty.

\subsection{Absence of analytical explanation}

All results in this paper are empirical. We have no analytical derivation
of the LIS ratio, no proof that it converges, and no theoretical
explanation for the specific value. The up-step fraction of $\sim 1/3$ is
well understood from the Collatz map's structure, but connecting this to a
precise LIS ratio requires understanding how LIS interacts with the
temporal distribution of up-steps, the peak position, and the suffix
length --- a problem we leave open.

% ======================================================================
\section{Conclusion}\label{sec:conclusion}

We have studied the longest increasing subsequence of Collatz trajectories
as a probe of pseudorandomness. The mean ratio $\LIS/2\sqrt{k}$ is
$0.567$ at $N = 10^6$, lying $0.22$ below the random permutation baseline
of $0.81$.

The primary contribution is not the headline deviation number but the
\emph{decomposition}. Through a hierarchy of null models, we show that:
\begin{itemize}
  \item The low ratio reflects temporal ordering (the shuffling test).
  \item The $31.6\%$ up-step fraction is the dominant driver (the Bernoulli
    model matches within $0.02$).
  \item The long descending suffix further depresses the ratio (prefix-only
    ratio is $0.85$, near baseline).
  \item The residual Collatz-specific signal is approximately $0.05$,
    small but reproducible.
\end{itemize}

The LIS ratio varies across the $3n+b$ family ($0.49$ to $0.61$),
suggesting it encodes dynamical information specific to each map variant.

Several questions remain open:
\begin{enumerate}
  \item Does $\bar{r}(N)$ converge as $N \to \infty$? If so, to what limit?
    Large-scale computation ($N \geq 10^8$) would help.
  \item Can the ratio be derived analytically from the known statistics of
    the Collatz map?
  \item Does the LIS ratio distinguish Collatz from other $3n+b$ maps in a
    dynamically meaningful way?
  \item Does the framework (LIS ratio as a pseudorandomness probe) yield
    similar insights for other deterministic integer sequences?
\end{enumerate}

Code and data reproducing all experiments are available at
\url{https://github.com/sevenevesai/CollatzLIS}.

% ======================================================================
\begin{thebibliography}{99}

\bibitem{BDJ1999}
J.~Baik, P.~Deift, and K.~Johansson,
On the distribution of the length of the longest increasing subsequence
of random permutations,
\emph{J.\ Amer.\ Math.\ Soc.}, 12 (1999), pp.~1119--1178.

\bibitem{AldousDiaconis1999}
D.~Aldous and P.~Diaconis,
Longest increasing subsequences: from patience sorting to the
Baik-Deift-Johansson theorem,
\emph{Bull.\ Amer.\ Math.\ Soc.}, 36 (1999), pp.~413--432.

\bibitem{TracyWidom1994}
C.~A.~Tracy and H.~Widom,
Level-spacing distributions and the Airy kernel,
\emph{Comm.\ Math.\ Phys.}, 159 (1994), pp.~151--174.

\bibitem{Tao2019}
T.~Tao,
Almost all orbits of the Collatz map attain almost bounded values,
\emph{Forum of Mathematics, Pi}, 10 (2022), e12.

\bibitem{Lagarias1985}
J.~C.~Lagarias,
The $3x+1$ problem and its generalizations,
\emph{Amer.\ Math.\ Monthly}, 92 (1985), pp.~3--23.

\bibitem{KontorovichMiller2005}
A.~V.~Kontorovich and S.~J.~Miller,
Benford's law, values of $L$-functions and the $3x+1$ problem,
\emph{Acta Arith.}, 120 (2005), pp.~269--297.

\bibitem{Terras1976}
R.~Terras,
A stopping time problem on the positive integers,
\emph{Acta Arith.}, 30 (1976), pp.~241--252.

\bibitem{Crandall1978}
R.~E.~Crandall,
On the ``$3x+1$'' problem,
\emph{Math.\ Comp.}, 32 (1978), pp.~1281--1292.

\bibitem{Cohen1988}
J.~Cohen,
\emph{Statistical Power Analysis for the Behavioral Sciences},
2nd ed., Lawrence Erlbaum Associates, 1988.

\bibitem{Efron1979}
B.~Efron,
Bootstrap methods: another look at the jackknife,
\emph{Ann.\ Statist.}, 7 (1979), pp.~1--26.

\end{thebibliography}

\end{document}
